\section{Docker Networks}

\begin{frame}
  \frametitle{Docker Networks}
  Controlam como containers se comunicam entre si e com o mundo externo.

  \vspace{0.5em}
  Três redes criadas automaticamente na instalação:

  \vspace{0.5em}
  \begin{center}
  \begin{tabular}{|l|l|l|}
  \hline
  \textbf{Nome} & \textbf{Driver} & \textbf{Descrição} \\
  \hline
  bridge & bridge & Rede virtual isolada (padrão) \\
  host & host & Compartilha rede do host \\
  none & null & Sem rede \\
  \hline
  \end{tabular}
  \end{center}
\end{frame}

\begin{frame}[fragile]
  \frametitle{Bridge (Padrão)}
  Cria uma rede virtual privada. Containers na mesma bridge se comunicam.

  \vspace{0.5em}
  \textbf{Bridge padrão}: comunicação apenas por IP (sem DNS).

  \vspace{0.5em}
  \textbf{Bridge customizada}: resolução DNS automática por nome!

  \vspace{0.5em}
\begin{verbatim}
docker network create minha-rede

docker run -d --name web --network minha-rede nginx
docker run -d --name api --network minha-rede myapp

# "api" pode acessar "web" pelo nome:
docker exec api curl http://web:80
\end{verbatim}
\end{frame}

\begin{frame}
  \frametitle{Host}
  Remove isolamento de rede. Container compartilha a pilha de rede do host.

  \vspace{0.5em}
  \texttt{docker run -d --network host nginx}

  \vspace{0.3em}
  Nginx acessível na porta 80 do host diretamente (sem \texttt{-p}).

  \vspace{0.8em}
  \begin{columns}
    \column{0.5\textwidth}
      \textbf{Quando usar:}
      \begin{itemize}
        \item Performance máxima
        \item Portas dinâmicas
        \item Monitoramento de rede
      \end{itemize}

    \column{0.5\textwidth}
      \textbf{Quando NÃO usar:}
      \begin{itemize}
        \item Isolamento necessário
        \item Mesma porta em vários containers
        \item Produção com segurança
      \end{itemize}
  \end{columns}

  \vspace{0.5em}
  \small \textbf{Obs}: Funciona nativamente apenas no Linux.
\end{frame}

\begin{frame}
  \frametitle{None}
  Desabilita completamente a rede do container.

  \vspace{0.5em}
  \texttt{docker run --network none alpine ping 8.8.8.8}

  \texttt{ping: sendto: Network unreachable}

  \vspace{0.8em}
  \begin{itemize}
    \item Apenas interface loopback (\texttt{127.0.0.1})
    \item Sem comunicação externa
    \item Útil para processamento local isolado
  \end{itemize}
\end{frame}

\begin{frame}
  \frametitle{Comandos de Rede}
  \small
  \texttt{docker network ls} \hspace{4em} — listar redes

  \vspace{0.2em}
  \texttt{docker network create <nome>} \hspace{0.5em} — criar rede

  \vspace{0.2em}
  \texttt{docker network inspect <nome>} — detalhes

  \vspace{0.2em}
  \texttt{docker network connect <rede> <cont.>} — conectar

  \vspace{0.2em}
  \texttt{docker network disconnect <rede> <cont.>} — desconectar

  \vspace{0.2em}
  \texttt{docker network rm <nome>} \hspace{1.5em} — remover

  \vspace{0.2em}
  \texttt{docker network prune} \hspace{3em} — remover não utilizadas
\end{frame}
